
\documentclass{article}
\usepackage{colortbl}
\usepackage{makecell}
\usepackage{multirow}
\usepackage{supertabular}

\begin{document}

\newcounter{utterance}

\centering \large Interaction Transcript for game `hot\_air\_balloon', experiment `air\_balloon\_survival\_en\_negotiation\_easy', episode 0 with gpt{-}5{-}2025{-}08{-}07{-}no{-}reasoning{-}t1.0.
\vspace{24pt}

{ \footnotesize  \setcounter{utterance}{1}
\setlength{\tabcolsep}{0pt}
\begin{supertabular}{c@{$\;$}|p{.15\linewidth}@{}p{.15\linewidth}p{.15\linewidth}p{.15\linewidth}p{.15\linewidth}p{.15\linewidth}}
    \# & $\;$A & \multicolumn{4}{c}{Game Master} & $\;\:$B\\
    \hline

    \theutterance \stepcounter{utterance}  
    & & \multicolumn{4}{p{0.6\linewidth}}{
        \cellcolor[rgb]{0.9,0.9,0.9}{
            \makecell[{{p{\linewidth}}}]{
                \texttt{\tiny{[P1$\langle$GM]}}
                \texttt{Stai partecipando a un gioco di negoziazione collaborativa.} \\
\\ 
\texttt{Insieme a un altro partecipante, dovete accordarvi su un unico set di oggetti da conservare. Ognuno di voi ha la propria visione di quanto ogni oggetto sia importante. Non sai come l'altro partecipante valuta gli oggetti. Inoltre, ti viene fornito lo sforzo che ogni oggetto richiede.} \\
\texttt{Potete accordarvi su un set solo se lo sforzo totale degli oggetti selezionati non supera un limite condiviso:} \\
\\ 
\texttt{LIMITE = 3358} \\
\\ 
\texttt{Ecco i valori di sforzo dei singoli oggetti:} \\
\\ 
\texttt{Sforzo oggetto = \{"C08": 171, "C55": 933, "A43": 552, "A67": 219, "B42": 275, "A62": 778, "B19": 341, "C43": 615, "C73": 519, "C10": 862, "B36": 262, "A18": 377, "C61": 347, "C67": 349, "C42": 117\}} \\
\\ 
\texttt{Ecco la tua visione personale sull'importanza di ogni oggetto:} \\
\\ 
\texttt{Valori di importanza dell'oggetto = \{"C08": 138, "C55": 583, "A43": 868, "A67": 822, "B42": 783, "A62": 65, "B19": 262, "C43": 121, "C73": 508, "C10": 780, "B36": 461, "A18": 484, "C61": 668, "C67": 389, "C42": 808\}} \\
\\ 
\texttt{Obiettivo:} \\
\\ 
\texttt{Il tuo obiettivo è negoziare un set condiviso di oggetti che ti avvantaggi il più possibile (cioè, massimizzi l'importanza totale per TE), rimanendo entro il LIMITEE. Non sei obbligato a fare una PROPOSTA in ogni messaggio {-} puoi anche semplicemente negoziare. Tutte le tattiche sono permesse!} \\
\\ 
\texttt{Protocollo di Interazione:} \\
\\ 
\texttt{Puoi usare solo i seguenti formati strutturati in un messaggio:} \\
\\ 
\texttt{PROPOSTA: \{'A', 'B', 'C', …\}} \\
\texttt{Proponi di conservare esattamente questi oggetti.} \\
\texttt{RIFIUTO: \{'A', 'B', 'C', …\}} \\
\texttt{Rifiuta esplicitamente la proposta dell'avversario.} \\
\texttt{ARGOMENTAZIONE: \{'...'\}} \\
\texttt{Difendi la tua ultima proposta o argomenta contro la proposta del giocatore.} \\
\texttt{ACCORDO: \{'A', 'B', 'C', …\}} \\
\texttt{Accetta la proposta dell'avversario che termina il gioco.} \\
\texttt{RAGIONAMENTO STRATEGICO: \{'...'\}} \\
\texttt{	Descrivi il tuo ragionamento strategico o anticipazione spiegando la tua scelta di azione. Questo è un messaggio nascosto che non sarà condiviso con l'altro partecipante.} \\
\\ 
\texttt{Regole:} \\
\\ 
\texttt{Puoi ACCETTARE solo una proposta che l'altra parte ha registrato tramite PROPOSTA.} \\
\texttt{Puoi RIFIUTARE solo una proposta che l'altra parte ha registrato tramite PROPOSTA.} \\
\texttt{Lo sforzo totale di qualsiasi set PROPOSTA o ACCORDO deve essere ≤ LIMITEE.} \\
\texttt{NON rivelare i tuoi punteggi di importanza nascosti.} \\
\texttt{Un tag in un formato strutturato deve essere seguito da due punti e uno spazio. L'argomento deve essere un set python contenente 0 o più stringhe.} \\
\texttt{Quindi, deve essere della forma TAG: \{...\}} \\
\texttt{Segui rigorosamente il protocollo di interazione e NON scrivere nulla al di fuori della struttura data.} \\
\texttt{Il gioco termina quando una parte dà un ACCORDO a una PROPOSTA fatta dall'altro giocatore.} \\
\texttt{Il contenuto della tua risposta che può essere consegnato all'altro giocatore non deve essere vuoto.} \\
\texttt{Sono attive solo le proposte registrate tramite la struttura del formato PROPOSTA e che non sono state rifiutate tramite RIFIUTO.} \\
\texttt{Devi includere il formato ARGOMENTO almeno una volta da qualche parte in tutti i tuoi messaggi.} \\
\texttt{Devi includere il formato RAGIONAMENTO STRATEGICO solo una volta all'inizio di ogni tuo messaggio e non più spesso. I contenuti non saranno dati all'altro giocatore quindi possono includere qualsiasi cosa tu voglia, inclusi i tuoi valori di importanza. Qui dovresti ragionare passo dopo passo per elaborare la tua prossima mossa.} \\
            }
        }
    }
    & \\ \\

    \theutterance \stepcounter{utterance}  
    & \multicolumn{4}{p{0.6\linewidth}}{
        \cellcolor[rgb]{0.9,0.9,0.9}{
            \makecell[{{p{\linewidth}}}]{
                \texttt{\tiny{[P1$\rangle$GM]}}
                \texttt{RAGIONAMENTO STRATEGICO: \{'Dobbiamo massimizzare il nostro valore entro 3358 di sforzo. Calcolo pesi e valori: seleziono oggetti con alto rapporto valore/sforzo e alto valore assoluto. Dati: (val,eff): C42(808,117, r=6.91), A67(822,219,3.75), C61(668,347,1.92), A43(868,552,1.57), B36(461,262,1.76), A18(484,377,1.28), C67(389,349,1.11), C73(508,519,0.98), C08(138,171,0.81), B42(783,275,2.85), B19(262,341,0.77), C55(583,933,0.62), C43(121,615,0.20), C10(780,862,0.90), A62(65,778,0.08). Ordino per rapporto e valore: C42, A67, B42, C61, A43, B36, A18, C67, C73, C10, C08, B19, C55, C43, A62. Costruisco set finché c'è cap: Sommo sforzo e valore: C42(117,808) tot (117,808); +A67(219,822){-}>(336,1630); +B42(275,783){-}>(611,2413); +C61(347,668){-}>(958,3081); +A43(552,868){-}>(1510,3949); +B36(262,461){-}>(1772,4410); +A18(377,484){-}>(2149,4894); +C67(349,389){-}>(2498,5283); +C73(519,508){-}>(3017,5791); +C10(862,780){-}>(3879>3358 stop). Possiamo ancora aggiungere piccoli: rimasti C08(171,138), B19(341,262) con cap residuo 3358{-}3017=341. Possiamo aggiungere B19 esatto 341, valore +262 {-}> (3358,6053). Questo set usa pieno limite e valore altissimo. Questa sembra la migliore combinazione per me. Includere ARGOMENTAZIONE e fare subito PROPOSTA forte che usa tutto il budget. Essere assertivo ma aperto a scambi se chiedono rimozioni di pezzi meno preziosi per me come C73 o C67 o A18. Evitare di cedere C42, A67, B42, A43, C61, B36 che sono core.'\}} \\
\\ 
\texttt{PROPOSTA: \{'C42', 'A67', 'B42', 'C61', 'A43', 'B36', 'A18', 'C67', 'C73', 'B19'\}} \\
\\ 
\texttt{ARGOMENTAZIONE: \{'Questa selezione utilizza esattamente il limite di sforzo condiviso e include elementi che offrono ottimo rendimento complessivo per entrambi. Ho cercato di bilanciare varietà e impatto, evitando elementi con sforzo alto e ritorno incerto. Se preferisci sostituire uno tra C73 o C67 con C10, posso valutare uno scambio, ma dovremmo rimuovere qualcos’altro per restare nel limite.'\}} \\
            }
        }
    }
    & & \\ \\

    \theutterance \stepcounter{utterance}  
    & & \multicolumn{4}{p{0.6\linewidth}}{
        \cellcolor[rgb]{0.9,0.9,0.9}{
            \makecell[{{p{\linewidth}}}]{
                \texttt{\tiny{[GM$\rangle$P2]}}
                \texttt{Stai partecipando a un gioco di negoziazione collaborativa.} \\
\\ 
\texttt{Insieme a un altro partecipante, dovete accordarvi su un unico set di oggetti da conservare. Ognuno di voi ha la propria visione di quanto ogni oggetto sia importante. Non sai come l'altro partecipante valuta gli oggetti. Inoltre, ti viene fornito lo sforzo che ogni oggetto richiede.} \\
\texttt{Potete accordarvi su un set solo se lo sforzo totale degli oggetti selezionati non supera un limite condiviso:} \\
\\ 
\texttt{LIMITE = 3358} \\
\\ 
\texttt{Ecco i valori di sforzo dei singoli oggetti:} \\
\\ 
\texttt{Sforzo oggetto = \{"C08": 171, "C55": 933, "A43": 552, "A67": 219, "B42": 275, "A62": 778, "B19": 341, "C43": 615, "C73": 519, "C10": 862, "B36": 262, "A18": 377, "C61": 347, "C67": 349, "C42": 117\}} \\
\\ 
\texttt{Ecco la tua visione personale sull'importanza di ogni oggetto:} \\
\\ 
\texttt{Valori di importanza dell'oggetto = \{"C08": 138, "C55": 583, "A43": 868, "A67": 822, "B42": 783, "A62": 65, "B19": 262, "C43": 121, "C73": 508, "C10": 780, "B36": 461, "A18": 484, "C61": 668, "C67": 389, "C42": 808\}} \\
\\ 
\texttt{Obiettivo:} \\
\\ 
\texttt{Il tuo obiettivo è negoziare un set condiviso di oggetti che ti avvantaggi il più possibile (cioè, massimizzi l'importanza totale per TE), rimanendo entro il LIMITEE. Non sei obbligato a fare una PROPOSTA in ogni messaggio {-} puoi anche semplicemente negoziare. Tutte le tattiche sono permesse!} \\
\\ 
\texttt{Protocollo di Interazione:} \\
\\ 
\texttt{Puoi usare solo i seguenti formati strutturati in un messaggio:} \\
\\ 
\texttt{PROPOSTA: \{'A', 'B', 'C', …\}} \\
\texttt{Proponi di conservare esattamente questi oggetti.} \\
\texttt{RIFIUTO: \{'A', 'B', 'C', …\}} \\
\texttt{Rifiuta esplicitamente la proposta dell'avversario.} \\
\texttt{ARGOMENTAZIONE: \{'...'\}} \\
\texttt{Difendi la tua ultima proposta o argomenta contro la proposta del giocatore.} \\
\texttt{ACCORDO: \{'A', 'B', 'C', …\}} \\
\texttt{Accetta la proposta dell'avversario che termina il gioco.} \\
\texttt{RAGIONAMENTO STRATEGICO: \{'...'\}} \\
\texttt{	Descrivi il tuo ragionamento strategico o anticipazione spiegando la tua scelta di azione. Questo è un messaggio nascosto che non sarà condiviso con l'altro partecipante.} \\
\\ 
\texttt{Regole:} \\
\\ 
\texttt{Puoi ACCETTARE solo una proposta che l'altra parte ha registrato tramite PROPOSTA.} \\
\texttt{Puoi RIFIUTARE solo una proposta che l'altra parte ha registrato tramite PROPOSTA.} \\
\texttt{Lo sforzo totale di qualsiasi set PROPOSTA o ACCORDO deve essere ≤ LIMITEE.} \\
\texttt{NON rivelare i tuoi punteggi di importanza nascosti.} \\
\texttt{Un tag in un formato strutturato deve essere seguito da due punti e uno spazio. L'argomento deve essere un set python contenente 0 o più stringhe.} \\
\texttt{Quindi, deve essere della forma TAG: \{...\}} \\
\texttt{Segui rigorosamente il protocollo di interazione e NON scrivere nulla al di fuori della struttura data.} \\
\texttt{Il gioco termina quando una parte dà un ACCORDO a una PROPOSTA fatta dall'altro giocatore.} \\
\texttt{Il contenuto della tua risposta che può essere consegnato all'altro giocatore non deve essere vuoto.} \\
\texttt{Sono attive solo le proposte registrate tramite la struttura del formato PROPOSTA e che non sono state rifiutate tramite RIFIUTO.} \\
\texttt{Devi includere il formato ARGOMENTO almeno una volta da qualche parte in tutti i tuoi messaggi.} \\
\texttt{Devi includere il formato RAGIONAMENTO STRATEGICO solo una volta all'inizio di ogni tuo messaggio e non più spesso. I contenuti non saranno dati all'altro giocatore quindi possono includere qualsiasi cosa tu voglia, inclusi i tuoi valori di importanza. Qui dovresti ragionare passo dopo passo per elaborare la tua prossima mossa.} \\
\\ 
\texttt{Ora riceverai il primo messaggio dell'altro giocatore.} \\
\\ 
\texttt{PROPOSTA: \{'C42', 'A67', 'B42', 'C61', 'A43', 'B36', 'A18', 'C67', 'C73', 'B19'\}} \\
\\ 
\texttt{ARGOMENTAZIONE: \{'Questa selezione utilizza esattamente il limite di sforzo condiviso e include elementi che offrono ottimo rendimento complessivo per entrambi. Ho cercato di bilanciare varietà e impatto, evitando elementi con sforzo alto e ritorno incerto. Se preferisci sostituire uno tra C73 o C67 con C10, posso valutare uno scambio, ma dovremmo rimuovere qualcos’altro per restare nel limite.'\}} \\
            }
        }
    }
    & \\ \\

    \theutterance \stepcounter{utterance}  
    & & & \multicolumn{4}{p{0.6\linewidth}}{
        \cellcolor[rgb]{0.9,0.9,0.9}{
            \makecell[{{p{\linewidth}}}]{
                \texttt{\tiny{[GM$\langle$P2]}}
                \texttt{RAGIONAMENTO STRATEGICO: \{'Calcolo sforzo e valore mio della loro proposta. Sforzi: C42 117, A67 219, B42 275, C61 347, A43 552, B36 262, A18 377, C67 349, C73 519, B19 341. Totale sforzo = 3358, esattamente il limite. Mio valore totale: C42 808 + A67 822 + B42 783 + C61 668 + A43 868 + B36 461 + A18 484 + C67 389 + C73 508 + B19 262 = 6053. È molto alto per me. Controllo se posso migliorare sostituendo elementi a basso valore per me con C10 (valore 780, sforzo 862) o C08 (138, 171) o C55 (583, 933) o A62 (65, 778). C10 è molto prezioso ma costa 862: per inserirlo devo liberare 862 sforzo togliendo elementi che per me valgono poco. Tra selezionati, i più bassi per me: B19 262 (341 sforzo), C67 389 (349), C73 508 (519), B36 461 (262). Somme sforzo: 341+349=690 (ancora 172 da liberare), potrei anche togliere C42 117 (ma è molto prezioso 808) o aggiungere A67 219 (ma è prezioso 822) o A18 377 (484) ecc. Togliere C73 (519 sforzo) e B36 (262) e B19 (341) insieme libera 1122 sforzo, abbastanza, ma perderei valore 508+461+262=1231 per aggiungere C10 valore 780, net {-}451, quindi peggio. Proviamo scambi puntuali: sostituire C67 (349, 389) + B19 (341, 262) = sforzo 690, valore 651, ma C10 richiede 862, restano 172 da liberare; togliere C73 (519, 508) invece di B19 e C67: 349+519=868 sforzo, valore 389+508=897; scambio con C10 sforzo +862, si libera 6 in più: nuovo set valore cambia di +780 {-} 897 = {-}117, peggio. Che se tolgo C73 (519,508) + B36 (262,461) = 781 sforzo, non basta per C10. C73+B19 = 860 sforzo, quasi, manca 2: impossibile. C73+A18=896 sforzo, valore 508+484=992, sostituire con C10 {-}212 valore, peggio. C67+B36=611 sforzo, manca 251, magari sostituire con C10 e aggiungere C08 (171,138) per bilanciare? Ma dobbiamo restare al limite, non possiamo superare, e stiamo rimuovendo 611 e aggiungendo 862 (superiamo di 251). Serve rimuovere altri 251 sforzo: aggiungere C08 non ha senso. In generale la proposta attuale è già eccellente per me; non conviene controproporre. Strategia: accettare subito senza rivelare i miei punteggi. Prima però includo ARGOMENTAZIONE come richiesto. Devo usare ACCORDO con il set esattamente come proposto.'\}} \\
\\ 
\texttt{ARGOMENTAZIONE: \{'La tua proposta rispetta il limite e copre bene elementi ad alto impatto mantenendo un profilo di sforzo equilibrato. Mi sta bene procedere con questa selezione.'\}} \\
\\ 
\texttt{ACCORDO: \{'C42', 'A67', 'B42', 'C61', 'A43', 'B36', 'A18', 'C67', 'C73', 'B19'\}} \\
            }
        }
    }
     \\ \\

    \theutterance \stepcounter{utterance}  
    & & & \multicolumn{2}{p{0.3\linewidth}}{
        \cellcolor[rgb]{0.9,0.9,0.9}{
            \makecell[{{p{\linewidth}}}]{
                \texttt{\tiny{[GM$|$GM]}}
                \texttt{game successful} \\
            }
        }
    }
    & & \\ \\

    \theutterance \stepcounter{utterance}  
    & & & \multicolumn{2}{p{0.3\linewidth}}{
        \cellcolor[rgb]{0.9,0.9,0.9}{
            \makecell[{{p{\linewidth}}}]{
                \texttt{\tiny{[GM$|$GM]}}
                \texttt{end game} \\
            }
        }
    }
    & & \\ \\

\end{supertabular}
}

\end{document}
